\manualchapter{DAC operation and routing}{sec:dac}
Disabling both Tone and Noise makes the channel act as a stepped DAC driven
by Level. In this mode the tuning controls instead offset and scale its CV:
\[
 V_\mathrm{effective}=(a+V_\mathrm{FM}/5)
 \bigl(V_\mathrm{Level}+(p+5)/2+V_\mathrm{pitch}/2\bigr),
\]
where $p$ is the frequency knob and $a$ is FM depth. The level knob multiplies
this result by $1/10$ before rounding/clamping to 0--15. Start with explicit
CV cables and modest offsets; the adjusted level voltage also forwards to
the next column. For ordinary pitched sound, turn Tone back on.

\subsection{Routing and troubleshooting}
The widest input sets 1--16 chip instances. Inputs use indexed channels;
match pitch and envelope channel counts. Each unpatched output forwards
its signal into the next, clipped at approximately $\pm5$\,V per stage.
Monitor C for a complete sum or patch earlier outputs separately.

Reset initializes the chip and envelope mode; phase is not stored. Tone
pitch, level and sync run each sample, while shared envelope/noise settings
and switches update every 16 samples. If a voice unexpectedly follows another
pitch, check normalling into its V/OCT and into Envelope V/OCT. If an envelope
runs only once, choose a repeating mode or provide sync triggers.
